\section{Regresión logística}
\label{sec:logistica}

\textbf{Diseño.} Para contrastar H4 modelamos la probabilidad de que un flujo sea ataque (\texttt{label} = 1) con regresión logística \citep{hosmer2013}. Dividimos la muestra en 3.500 flujos de entrenamiento (51,8\,\% ataques) y 1.500 de prueba (51,9\,\%), estratificando por \texttt{label} (semilla 42) (no usamos el testing set original: no está en el repositorio y habría que prepararlo igual). Para evitar fuga reajustamos el ACP solo con el entrenamiento y proyectamos la prueba; las cargas son casi idénticas a las del capítulo~\ref{sec:acp} (correlación $\ge 0{,}999$). Los predictores son PC1 a PC3 y \texttt{proto\_g} (referencia tcp), \texttt{service\_g} (ninguno) y \texttt{state\_g} (FIN), con la categoría más frecuente como referencia. Excluimos TTL (separación perfecta artefactual) y \texttt{tipo\_trafico} y \texttt{attack\_cat} (fuga de la etiqueta).

\textbf{Elección del modelo.} M1 usa solo los puntajes; M2 agrega protocolo, servicio y estado. La razón de verosimilitud favorece a M2 ($\chi^2 = 862{,}2$; 7 gl; $p < 0{,}001$), con unos 850 puntos menos de AIC y 0,08 más de AUC (Tabla~\ref{tab:log-modelos}). Elegimos M2.

\begin{table}[H]
  \centering
  \caption{Comparación de modelos; el AUC es sobre la muestra de prueba.}\label{tab:log-modelos}
  \begin{tablacompacta}
  \begin{tabular}{lrrr} \toprule
  Modelo & AIC & $R^2$ McFadden & AUC prueba \\ \midrule
  M1: 3 puntajes del ACP & 3.744,1 & 0,229 & 0,794 \\
  M2: M1 + protocolo, servicio, estado & 2.895,9 & 0,407 & 0,878 \\ \bottomrule
  \end{tabular}
  \end{tablacompacta}
\end{table}

\textbf{Odds ratios.} Con lo demás fijo, un desvío más de PC1 (más volumen y temporización, menos repetición) multiplica las chances de ataque por 4,7, y uno más de PC2 (más volumen y repetición a la vez) las reduce un 39\,\% (Tabla~\ref{tab:log-or}). El estado CON\_otros tiene chances mucho menores que FIN (OR 0,05). Los OR de las categóricas son condicionales y poco estables por colinealidad (VIF $\approx 16$ en \texttt{udp} e INT; $\le 3{,}7$ en las componentes), así que no deben leerse de a uno: dns tiene 31\,\% de ataques en el entrenamiento (49\,\% en servicio ninguno) pero OR ajustado 18,7: el signo cambia al controlar por estado y protocolo. El IC de INT es muy ancho: no debe citarse ``93 veces''.

\begin{table}[H]
  \centering
  \caption{Odds ratios de M2 con IC 95\,\% por verosimilitud perfilada (componentes: por desvío del puntaje; categóricas: frente a la referencia). Los términos de protocolo (udp, otros) y servicio otros tienen IC que incluyen 1.}\label{tab:log-or}
  \begin{tablacompacta}
  \begin{tabular}{lrr} \toprule
  Término & OR & IC 95\,\% \\ \midrule
  PC1 & 4,68 & 3,82 -- 5,77 \\
  PC2 & 0,61 & 0,54 -- 0,69 \\
  PC3 & 2,55 & 2,27 -- 2,89 \\
  Estado INT vs FIN & 92,7 & 22,6 -- 361,8 \\
  Estado CON\_otros vs FIN & 0,05 & 0,02 -- 0,13 \\
  Servicio dns vs ninguno & 18,7 & 8,5 -- 43,0 \\
  Servicio http vs ninguno & 1,88 & 1,47 -- 2,42 \\ \bottomrule
  \end{tabular}
  \end{tablacompacta}
\end{table}

\textbf{Desempeño en prueba.} Con umbral 0,5, M2 logra AUC = 0,878 (IC 95\,\% de DeLong con \texttt{pROC} \citep{robin2011}: 0,860 -- 0,896), exactitud 0,817, sensibilidad 0,910, especificidad 0,717 y precisión 0,776. Domina el falso positivo: 204 normales marcados como ataque (28,3\,\% de los normales) frente a 70 ataques no detectados (9,0\,\%; Figura~\ref{fig:log-roc}). En detección de intrusiones el modelo prioriza no dejar pasar ataques a costa de alertas falsas.

\begin{figure}[H]
  \centering
  \grafico[0.6\linewidth]{17_logistica_roc_confusion.png}
  \caption{Curva ROC y matriz de confusión de M2 en la prueba (umbral 0,5). La curva se aleja de la diagonal del azar; el error principal son los normales clasificados como ataque. Fuente: elaboración propia con UNSW-NB15.}
  \label{fig:log-roc}
\end{figure}

\textbf{Diagnóstico.} No hay separación: convergió en 6 iteraciones sin avisos, el coeficiente máximo es 4,5 y ninguna probabilidad ajustada está a menos de $10^{-8}$ de 0 o 1. En influencia, 151 flujos (4,3\,\%) superan $4/n$ en la distancia de Cook, pero el máximo es 0,030, muy lejos de 1. La calibración es razonable, salvo que el decil 10 sobreestima (predicho 0,94, observado 0,86; Anexo A, Figura~\ref{fig:a-log-or-cal}).

\textbf{Cierre.} H4 se confirma: el AUC en prueba es 0,878 y el límite inferior de su IC (0,860) también supera 0,80; con las componentes solas (M1, 0,794) no se cumpliría. La probabilidad estimada es relativa a esta base (52\,\% de ataques, sin duplicados, ataques sintéticos): interpretamos odds ratios, no prevalencias, y no extrapolamos a una red real.
